\documentclass[acmsmall,anonymous]{acmart}

\usepackage{microtype}
\usepackage{enumitem}
\usepackage{amsthm}
\usepackage{caption}
\usepackage{booktabs}


% ---------- Notation ----------
\newcommand{\Hist}{\mathcal{H}}
\newcommand{\Spec}{\mathsf{Spec}}
\newcommand{\drule}{\mathrel{:\!-}}
\newcommand{\Prov}{\mathsf{Prov}}
\newcommand{\snap}{\mathsf{snap}}
\newcommand{\seal}{\mathsf{seal}}
\newcommand{\FCW}{\mathsf{FCW}}
\newcommand{\seq}{\mathbin{\triangleright}}

% ---------- Theorem environments ----------
\newtheorem{definition}{Definition}
\newtheorem{theorem}{Theorem}
\newtheorem{lemma}{Lemma}
\newtheorem{corollary}{Corollary}
\theoremstyle{remark}
\newtheorem{remark}{Remark}
\newtheorem{example}{Example}
\newtheorem{observation}{Observation}
\newtheorem{assumption}{Assumption}

\numberwithin{definition}{section}
\numberwithin{theorem}{section}
\numberwithin{lemma}{section}
\numberwithin{corollary}{section}
\numberwithin{remark}{section}
\numberwithin{example}{section}
\numberwithin{observation}{section}
\numberwithin{assumption}{section}

\title{Provenance of Determinations: From Ambiguity to Algebra}
\author{Joseph M. Hellerstein}
\affiliation{%
  \institution{University of California, Berkeley and Amazon Web Services}
  \country{USA}
}
\date{\today}

\acmConference[Conference]{Full Conference Name}{Month Year}{Location}
\acmBooktitle{Proceedings of the Conference}
\acmYear{2026}
\acmISBN{123-4-56789-012-3}


\begin{document}

\begin{abstract}
  Algebraic provenance explains query results by tracing derivations under a fixed,
  resolved semantics.
  Many systems, however, operate under \emph{semantic ambiguity}:
  logic programs with negation admit multiple accepted models, transactional
  systems admit multiple consistent executions, and distributed or policy-aware
  systems admit multiple outcomes compatible with the same history.
  In such settings, it is not enough to explain \emph{what follows;} one must also
  explain \emph{under which semantic resolution of ambiguity} it follows.

  This paper introduces \emph{determination provenance}, a framework for making
  semantic resolution explicit and explainable.
  A \emph{determination} is a structured collection of commitments that
  resolves ambiguity by refining admissible outcomes.
  We show that provenance becomes meaningful only \emph{after} such resolution.
  % , and
  % that classical algebraic provenance arises as a special case when semantics are
  % already resolved.
  We formalize how determinations structure provenance prior to resolution, giving
  rise to a layered, dependency-sensitive form that becomes ordinary semiring
  provenance once a resolving determination is fixed.
  We introduce \emph{layer certificates} as sufficient summaries of
  commitments for downstream reasoning.

  The work extends the reach of formal data provenance to a range of
  settings where ambiguity is intrinsic.
  To illustrate, we instantiate the framework in two domains with fundamentally different sources
  of ambiguity: negation in Datalog and scheduling in transactional systems.
  Across both settings, determination provenance clarifies which
  commitments were required to justify observations, and why determinations
  are required to complement classical provenance in the presence of ambiguity.
\end{abstract}

\maketitle

% ================================================================
\section{Introduction}
\label{sec:intro}

Algebraic provenance has given database theory a powerful and durable language for
explaining query results.
Declarative query languages provide well-defined semantics,
and provenance provides clean, uniform explanations of query
results—supporting why- and how-provenance, uncertainty propagation,
and optimization through a common algebraic framework
\cite{green2007provenance,suciu2011probabilistic}.
In these settings, explanation reduces to arithmetic, and this reduction has been
remarkably effective.

Practical data systems can make use of this foundation, but they operate in a
broader semantic landscape.
Transactional workloads admit multiple serializable executions.
Logic programs with negation admit multiple distinct models.
Distributed systems can yield multiple admissible orderings consistent with the same
history.
In such settings, a history determines not a single outcome, but a structured
space of multiple admissible outcomes.

Natural provenance questions often require more involved answers in this context
than classical provenance can provide.
For example, the reason that a tuple was
returned by a query may depend on which serializable execution was realized by a scheduler.
New questions arise in this context as well that have no analog in classical provenance.
Did this tuple appear as an inevitable consequence regardless of transactional decisions,
or because of a particular scheduling choice made by the system?
Was this tuple returned due to a particular choice of isolation
level---e.g., snapshot isolation versus serializability?
These familiar transactional examples are emblematic of the gap between classical data provenance and a range of
systems research on
lineage~\cite{sigelman2010dapper,seltzer2005provenance,bates2015trustworthy}
that has lacked a formal foundation.
Answering such questions requires explaining not only how conclusions were derived,
but \emph{how semantic ambiguity was resolved}.

Prior work developed
\emph{determinations} as formal objects that resolve semantic ambiguity by
introducing structured commitments~\cite{hellerstein2026interpretation}.
That work develops a complexity theory of determination, characterizing the
\emph{cost} and \emph{depth} of semantic resolution.
Determinations explain both semantic ambiguity arising within formal languages
(as in logic programming with negation) and ambiguity arising from execution
environments (as in transactional and distributed systems).

This paper shows how determinations give rise to a new form of provenance.
We introduce \emph{determination provenance}, which records the
commitments under which observations become well defined.
Once semantic ambiguity has been resolved by a determination, classical data
provenance applies directly, providing the derivational structure of results
within that determination.

We illustrate determination provenance in two domains with fundamentally different
sources of semantic ambiguity.
Logic programming with negation provides the canonical, mathematically explicit
instance of semantic ambiguity and resolution in computer science, serving as a
control experiment for determination provenance.
We then turn to transactional systems, where ambiguity arises from concurrency and
isolation choices, and show how determination provenance provides a fuller
accounting of query results than classical approaches. In the appendix we also examine
examples from distributed systems observability, a widely-used setting for system
lineage~\cite{sigelman2010dapper}.
Here we give an initial transactional example to
illustrate the ideas:

\begin{example}[Provenance under transactional ambiguity]
  \label{ex:example}
  Consider a database with a binary relation $S$
  and an initially empty instance (EDB).
  The following transactions execute concurrently.
  % Ordering between two transactions is determined by the mutual ordering of their
  % respective commit and begin events.
  \[
    \begin{array}{ll}
      T_1: \textsf{begin};\ \textsf{insert } S(1,b)\ (x_1);\ \textsf{commit} \qquad &
      T_2: \textsf{begin};\ \textsf{insert } S(2,b)\ (x_2);\ \textsf{commit}          \\
      T_3: \textsf{begin};\ \textsf{insert } S(3,c)\ (x_3);\ \textsf{commit} \qquad &
      T_4: \textsf{begin};\ \textsf{insert } S(4,d)\ (x_4);\ \textsf{commit}          \\
      T_5: \textsf{begin};\ \textsf{delete } S(4,d);\ \textsf{commit}.              &
    \end{array}
  \]
  Subsequently, a query transaction $T_Q$ executes
  $\textsf{begin};\ Q(y) \drule S(k,y);\ \textsf{commit}$.

  Under serializability, the observed history admits multiple executions.
  Tuples $\{(b),(c)\}$ appear in $Q$ for every serializable execution.
  Tuple $(d)$ is contingent: it appears in some executions and not others, depending
  on whether $T_5$ commits before or after $T_Q$.
  Before resolving this ambiguity, even the \emph{family} of provenance expressions
  for $(d)$ is ambiguous: in one resolution the relevant question is \emph{why} $d$
  appears, while in another it is \emph{why-not} $d$.
  Determination provenance makes this dependence explicit by recording the
  commitments under which the query result is determined, together with the ordinary
  derivational explanation that applies once those commitments are fixed.
  Section~\ref{sec:transactions} formalizes this example by writing down the
  corresponding resolving determinations and provenance expressions, and the example
  is worked in full in Appendix~\ref{sec:appendix-concurrency}.
\end{example}

\paragraph{Contributions.}
This paper introduces \emph{determination provenance}: a semantic layer that
explains conclusions in the presence of ambiguous semantics by recording the
commitments under which those conclusions become determinate.
We (i) formalize determinations and certificates as the semantic structure
needed to recover classical algebraic provenance once ambiguity is resolved;
(ii) show, via Datalog with negation, how different semantic regimes induce
distinct determination and provenance structures, (iii) show, via transactional
systems, how concurrency and isolation choices give rise to determination-dependent
query results; and (iv) identify \emph{robustness} across determinations and
\emph{preservation} across changes in specifications
as new classes of provenance questions enabled by this framework.


% ================================================================
\section{Preliminaries: Ambiguous Semantics and Refinement}
\label{sec:prelim}
% ================================================================

Classical data provenance assumes a fixed, resolved semantics: for each input instance,
a specification (query) determines a single outcome, and provenance explains how that outcome was
derived.
In many systems, however, concurrency, negation, or underspecification admit
multiple admissible outcomes for the same input.

This section introduces a minimal semantic framework for reasoning about such
ambiguity.
We model specifications with multiple admissible outcomes and the
commitments—called \emph{determinations}—that restrict ambiguity until meaning
is determinate.
Our focus is purely structural: we introduce only the semantic objects needed to
define determination provenance and to reason about robustness and change across
determinations, without assuming any particular execution model.

\subsection{Histories}

When semantics are ambiguous, it is no longer sufficient to model input as a
single static instance.
Instead, we require a representation of \emph{partial information}: what has
been observed, and which semantic distinctions remain unresolved.

We model such partial information using \emph{histories}.
A history records a collection of semantic events together with any known
precedence constraints among them.
Histories generalize familiar notions of input:
in logic programming, a history may consist of a set of ground facts;
in transactional systems, it may record transaction invocations, reads, writes,
and commits; in distributed settings, it may include message deliveries or
observations.
When no precedence constraints are present, histories collapse to unordered
sets and coincide with classical database instances.

\begin{definition}[History]
  A \emph{history} is a finite partially ordered set
  \[
    H = (E, \rightarrow),
  \]
  where $E$ is a set of events and $\rightarrow \subseteq E \times E$ is a strict
  (irreflexive and transitive) partial order representing known precedence constraints among events.
\end{definition}

The order $\rightarrow$ records only irreducible constraints—facts that must be
respected by any admissible semantic determination.
Histories abstract semantic information from execution detail: they may coincide
with executions or traces, but they do not require a total order or a fixed
schedule, and apply equally to purely semantic settings such as Datalog with
negation.

We write $\Hist$ for the class of histories under consideration.
This choice is parametric and deliberately abstract: histories may represent
database updates, transaction invocations, message deliveries, or any other
events that constrain semantic determination.

\begin{definition}[History Extension]
  For histories $H_1=(E_1,\rightarrow_1)$ and $H_2=(E_2,\rightarrow_2)$, we write
  $H_1 \sqsubseteq H_2$ if:
  \begin{enumerate}
    \item $E_1 \subseteq E_2$, and
    \item $\rightarrow_1 = \rightarrow_2 \cap (E_1 \times E_1)$, and
    \item no event in $E_1$ has a causal predecessor in $E_2 \setminus E_1$.
  \end{enumerate}
\end{definition}

Intuitively, a history extension may add new events and causal edges involving
those new events, but it may not revise the causal past of any event that has
already occurred.
We interpret $\sqsubseteq$ as information growth, not as time or execution \emph{per se}.

\subsection{Observations and Specifications}
Fix a set $O$ of fully specified semantic objects, called \emph{outcomes}.
\begin{definition}[Observation]
  An \emph{observation} is an element of a set $O$ of fully specified semantic objects.
\end{definition}

Observations represent what an observer may conclude once all semantic commitments have
been fixed.
Examples include a database state, a serial execution of transactions, a model of
a logic program, the visible state of an automaton, or a completed task result.

\begin{definition}[Specification]
  A \emph{specification} is a function
  \[
    \Spec : \Hist \to 2^{O},
  \]
  mapping each history $H \in \Hist$ to a (possibly empty) set of admissible observations.
\end{definition}

The set $\Spec(H)$ captures all observations consistent with the evidence in $H$.
When $|\Spec(H)| > 1$, the specification admits multiple admissible meanings for $H$,
reflecting unresolved semantic ambiguity.
When $|\Spec(H)| = 1$, the meaning of $H$ is fully determined.

Ambiguity in this sense is common and intentional.
Specifications for Datalog with negation, transactional systems, and distributed
executions all admit multiple admissible observations for the same history,
corresponding to different semantic resolutions of underspecified behavior.
In such settings, semantic questions cannot be posed until ambiguity is constrained
by determination.\footnote{%
  The ambiguity we study lives in the \emph{powerset domain} $2^{O}$: multiplicity
  of admissible outcomes $\Spec(H) \subseteq O$.
  This is distinct from uncertainty or nondeterminism \emph{within} the outcome
  domain $O$ itself (e.g., individual final outcomes in the form of distributions or sets of possible worlds).
  This distinction appears in more detail in Section~\ref{sec:negation} and
  Appendix~\ref{sec:appendix-datalog}.}


\subsection{Semantic Refinement}

Ambiguity is constrained by making additional semantic commitments.
We model these commitments as semantic conditions that \emph{filter} the set of admissible observations while
preserving feasibility.

\begin{definition}[Commitment (as Outcome Filter)]
  A \emph{commitment} is a function
  \[
    \varphi : \Hist \times 2^{O} \to 2^{O}
  \]
  such that for all histories $H$ and observation sets $S \subseteq O$, \hspace{0.8mm} $\varphi(H,S)
    \subseteq S$ (\emph{Shrinkage}), and if $S \neq \emptyset$ then $\varphi(H,S) \neq \emptyset$
  (\emph{Feasibility}).
\end{definition}
Commitments restrict admissible observations without introducing new ones or rendering
the specification infeasible.
Note that the Feasibility condition is specific to the provenance setting:
the companion complexity theory~\cite{hellerstein2026interpretation} requires only
Shrinkage, because infeasible commitments can arise naturally in adversarial or
exploratory resolution strategies.
Here, we require Feasibility because provenance is meaningful only when the
conditioned specification remains non-empty.

Crucially, a commitment may depend on the admissible outcome set
\emph{as a whole}.
That is, $\varphi(H,S)$ need not be definable as a pointwise test applied
independently to each $o \in S$, but may instead enforce semantic commitments
whose effect depends on which alternatives are jointly admissible.
We refer to such commitments as \emph{ambiguity-sensitive}.

\begin{definition}[Refinement by Conditioning]
  Given a specification $\Spec$ and a commitment $\varphi$, the
  \emph{refinement} of $\Spec$ by $\varphi$ is the specification
  \[
    (\Spec \mid \varphi)(H) \triangleq \varphi(H,\Spec(H)).
  \]
\end{definition}

Multiple commitments may be applied by iterating this refinement
operation.
Because commitments may be ambiguity-sensitive, these refinements \emph{need not commute}.

\subsection{Determinations}

Individual commitments typically express only partial semantic
commitments.
We therefore model semantic resolution using finite compositions of
commitments, applied as \emph{sequential conditioning}.

\begin{definition}[Determination]
  A \emph{determination} is a finite expression
  \[
    I = \varphi_1 \wedge \cdots \wedge \varphi_m
  \]
  of commitments, where $\wedge$ denotes \emph{ordered (sequential)
    refinement} rather than logical conjunction.
  In particular, $\wedge$ is not assumed to be commutative or idempotent.
\end{definition}

Refinement by a determination is defined by sequential application:
\[
  \Spec \mid I
  \;\triangleq\;
  (\cdots((\Spec \mid \varphi_1)\mid \varphi_2)\cdots)\mid \varphi_m .
\]

Because commitments may be ambiguity-sensitive, the order of conjuncts
can matter.
This non-commutativity induces semantic dependency.

\paragraph{Layer-sequencing notation.}
The operator $\wedge$ denotes \emph{sequential refinement} and is
therefore generally non-commutative.
When reasoning about \emph{depth}, however, we will often group commitments into
\emph{layers} whose commitments commute (i.e., can be applied in any order without
changing the result).
To make this structure explicit, we use $\seq$ as syntactic sugar for a sequence
of commuting batches.

\begin{definition}[Sequencing of layers (notation)]
  \label{def:layer-seq}
  Let $L_1,\ldots,L_k$ be multisets of commitments.
  We write $L_1 \seq L_2 \seq \cdots \seq L_k$
  for the determination obtained by applying the commitments of $L_1$ (in any
  order), then those of $L_2$, and so on:
  \[
    \Spec \mid (L_1 \seq \cdots \seq L_k)
    \;\triangleq\;
    (\cdots((\Spec \mid L_1)\mid L_2)\cdots)\mid L_k,
  \]
  where $\Spec \mid L$ abbreviates refinement by all commitments in $L$ in any
  order.
  This notation is used only to expose layering; it imposes no additional algebraic
  structure beyond commutativity within each layer.
\end{definition}

\begin{definition}[Resolved specification and resolving determination]
  A specification $\Spec$ is \emph{resolved} if for every history $H$, $|\Spec(H)| = 1$.
  A determination $I$ \emph{resolves} $\Spec$ (i.e., $I$ is \emph{resolving for} $\Spec$) if $\Spec \mid I$ is resolved.
\end{definition}

\paragraph{Clarification.}
Resolution is a property of the set-valued specification $\Spec:\Hist\to 2^O$:
a resolved specification
is \emph{single-valued pointwise in the history}, i.e., for every history $H$,
the set $\Spec(H)$ is a singleton.
Resolution marks the point at which semantic ambiguity has been fully discharged.

\subsection{Restriction and Identification}

Fix a specification $\Spec$ and a determination $I$.
For each history $H$, refinement by $I$ produces a restricted set of admissible
outcomes $(\Spec \mid I)(H)$.
The semantic effect of a determination has two aspects:
\emph{restriction}, in which outcomes incompatible with the commitments in $I$
are eliminated, and \emph{identification}, in which remaining outcomes that no
commitment in $I$ distinguishes become semantically equivalent.

Formally, write $O_H^I \triangleq (\Spec \mid I)(H)$ and define
$o \sim_{I,H} o'$ iff no commitment in $I$ distinguishes $o$ from $o'$ at $H$.
This equivalence induces a quotient $O_H^I / {\sim_{I,H}}$ for each history,
recording which admissible outcomes remain semantically distinguished after
conditioning.
Because commitments may depend on the history, this quotient is
history-indexed rather than global.
The quotient structure underlies the parsimony analysis in
Section~\ref{sec:parsimony}.

\subsection{Resolution Induces Monotonicity}
\label{sec:mono-theorem}

The central claim of this paper is that classical provenance applies once a
determination has resolved a specification.
Classical provenance requires monotone evaluation: facts accumulate and
explanations only grow.
We now show that monotonicity is not an independent assumption but a
\emph{consequence} of resolution, under a mild condition on the commitment basis.

\begin{definition}[History-monotone specification]
  \label{def:set-mono}
  A specification $\Spec : \Hist \to 2^O$ is \emph{history-monotone} if
  $H_1 \sqsubseteq H_2 \;\Longrightarrow\; \Spec(H_2) \subseteq \Spec(H_1)$.
\end{definition}

Extending a history may rule out outcomes but should never make new ones
admissible.

\begin{definition}[Persistent commitment basis]
  \label{def:persistent}
  A commitment basis $\Phi$ is \emph{persistent} if for every
  $\varphi \in \Phi$, all histories $H_1 \sqsubseteq H_2$,
  all $S, S' \subseteq O$, and all outcomes $o \in O$,
  \[
    o \notin \varphi(H_1, S)
    \quad\Longrightarrow\quad
    o \notin \varphi(H_2, S').
  \]
  That is, once a commitment excludes an outcome at a history,
  it continues to exclude that outcome at all extensions and under
  any admissible set.
\end{definition}

Persistence is a natural condition: it says that commitments are irrevocable.
The atomic basis (single-outcome exclusions) is trivially persistent.
The sealing basis used for Datalog with negation and the ordering basis used for
serializability are both persistent, as are snapshot predicates under SI.

\begin{assumption}[Prefix-closure of histories]
  \label{ass:prefix-closed}
  If $H \in \Hist$ and $H' \sqsubseteq H$, then $H' \in \Hist$.
\end{assumption}

\begin{theorem}[Resolved specifications are history-monotone]
  \label{thm:resolution-implies-monotone}
  Assume prefix-closure (Assumption~\ref{ass:prefix-closed}).
  Let $\Phi$ be a persistent commitment basis and let $I$ be a determination
  over $\Phi$ that resolves $\Spec$.
  Then the induced singleton-valued specification
  $\widehat{\Spec}_I(H) \triangleq (\Spec \mid I)(H)$
  is history-monotone.
\end{theorem}

\begin{proof}
  Fix $H_1 \sqsubseteq H_2$.
  Since $\Spec \mid I$ is resolved,
  $(\Spec \mid I)(H_1) = \{o_1\}$ and
  $(\Spec \mid I)(H_2) = \{o_2\}$ for some $o_1, o_2 \in O$.
  We show $o_1 = o_2$.

  Write $I = \varphi_1 \wedge \cdots \wedge \varphi_m$.
  Since $I$ refines $\Spec(H_1)$ to $\{o_1\}$, every outcome
  $o \neq o_1$ is excluded by some $\varphi_j$ when applied at $H_1$:
  there exists $j$ and an admissible set $S$ such that
  $o \notin \varphi_j(H_1, S)$.
  By persistence (Definition~\ref{def:persistent}), $o \notin \varphi_j(H_2, S')$
  for \emph{every} $S' \subseteq O$—in particular, for the admissible set that
  arises when $\varphi_j$ is applied at $H_2$ during the sequential refinement.
  Hence $o \notin (\Spec \mid I)(H_2)$.
  Since this holds for every $o \neq o_1$, the only outcome that can survive
  at $H_2$ is $o_1$, so $o_2 = o_1$.
\end{proof}

\begin{remark}
  The theorem does not claim that $\Spec$ itself is monotone, nor that all
  determinations coincide.
  It shows only that resolution forces monotonicity: once a determination fixes
  a single outcome at each history, the induced semantics must respect causal
  extension.
  Persistence should be verified for each application-specific basis; the
  examples in this paper all satisfy it.
\end{remark}

These notions of refinement, restriction, identification, and the monotonicity
of resolved specifications provide the semantic substrate on which we build
provenance structure in the next section.

%% ================================================================
\section{Determinations and Provenance Structure}
\label{sec:conditional-semiring}
This section connects the semantic notion of determination
(Section~\ref{sec:prelim}) to provenance.
Fix a specification $\Spec$ and a resolving determination $I$.
Determination provenance records both the semantic commitments fixed by $I$ and
the derivational provenance that applies after conditioning on $I$.

\subsection{Syntax of Determination Provenance}

We now introduce a compact notation for \emph{determination provenance}.
This subsection introduces no new semantic machinery.
It packages together notions already defined—determinations, resolution, and
classical provenance—into a single explanatory object.
Throughout the paper, determination provenance is expressed by conditioning
semantic objects on a determination.
We assume no algebraic structure on determinations or outcomes
beyond what has already been defined.

\paragraph{Determination provenance values.}
Fix a commutative semiring $(K,\otimes,1)$ for derivational provenance
(e.g., polynomial provenance).
For a query $Q$ and tuple $t$, write $\Prov_{Q \mid I}(t) \in K$ for the
ordinary (depth--0) provenance computed under the conditioned specification
$\Spec \mid I$.
We define \emph{determination provenance} as the pair
\[
  \mathsf{DP}_Q(t) \;\triangleq\; (I,\, \Prov_{Q \mid I}(t)).
\]

% When convenient, we may write $I \cdot k$ as an abbreviation for the pair $(I,k)$,
% but all semantics are expressed via conditioning on $I$.

The intended reading is:
\emph{$k$ explains the derivation, assuming the semantic commitments fixed by $I$.}

\paragraph{Relationship to classical provenance.}
Once a determination $I$ resolves $\Spec$, semantic ambiguity is fixed for the purpose of provenance.
Query evaluation under $\Spec \mid I$ is single-valued, and classical semiring
provenance applies without modification.
Determination provenance therefore records classical provenance under
conditioning, together with the determination that justifies that conditioning.

In particular, for \emph{depth--0} specifications—where no nontrivial semantic
commitments are required—determination provenance coincides exactly with
classical semiring-based provenance.

\paragraph{Why determination provenance is determination-relative.}
Determination provenance is defined relative to a \emph{single, fixed}
determination.
A provenance value records how a semantic outcome is derived \emph{given} a
particular set of semantic commitments; it does not combine or quantify over
multiple determinations.

Accordingly, we do not form disjunctions or aggregations over determinations
within a single provenance object.
Questions about how outcomes or explanations vary across different
determinations—such as robustness, invariance, or sensitivity to semantic
change—are treated explicitly and separately in
Section~\ref{sec:robustness}.

We illustrate this syntax next in a transactional setting, where
different ordering commitments induce different determinations and,
consequently, different provenance for the same query.

\subsection{Algebraic Structure of Determinations}
\label{sec:algebra}

Determinations refine a specification through a sequence of semantic
commitments.
This process exhibits a layered algebraic structure, which we summarize here to
clarify the relationship between determination provenance and classical algebraic
provenance.

\paragraph{Sequential conditioning across layers.}
Between layers, commitments compose by sequential conditioning.
This composition is generally dependent and non-commutative: applying one
commitment may enable or disable another.
Accordingly, the sequencing operator $\seq$ carries no algebraic meaning beyond
ordered refinement.

\paragraph{Commutative structure within a layer.}
Within a single layer, commitments commute.
Order of application does not affect the resulting restriction of admissible
outcomes.
As a result, commitments within a layer form a \emph{conditional commutative
  monoid}, where conjunction corresponds to joint semantic commitment and the
identity element is the trivial commitment.

The carrier of this monoid is conditional on prior layers: which predicates are
available, and which outcomes they act upon, depends on the determination prefix.

\paragraph{Special case: semiring structure.}
In special cases—most notably when commitments are pointwise and
independent across outcomes—this monoid extends to a commutative semiring.
Although this structure is familiar from algebraic provenance, we do not exploit
it here, as our focus is on semantic resolution rather than algebraic
optimization.

\paragraph{Why semiring provenance applies at layer~0.}
Once a resolving determination $I$ is fixed, all semantic ambiguity has been
eliminated.
The conditioned specification $\Spec \mid I$ is single-valued and, under a
persistent basis, history-monotone
(Theorem~\ref{thm:resolution-implies-monotone}): no
further commitments can rule out admissible outcomes.
All remaining variation is purely derivational.
In this setting, explanations combine algebraically exactly as in classical
semiring provenance.
While $\Spec \mid I$ need not syntactically be a Datalog program, it is
algebraically equivalent, for provenance purposes, to monotone query evaluation
over a fixed instance.
This is precisely the semantic regime studied in prior work on semiring
provenance.

Accordingly, determination provenance pairs a resolving determination with
ordinary algebraic provenance computed under the resolved specification.
Formal details appear in Appendix~\ref{app:algebra}.

\subsection{Example Revisited: Resolving Transactional Ambiguity}
\label{sec:concurrency}

We now return to the transactional Example~\ref{ex:example}.
Recall that values $b, c$ are present in every serializable execution, while value $d$
is present only under certain execution orders.

We can now show the syntax of determination provenance for this example.
In a conflict serializability context, commitments correspond to ordering
commitments among transactions: $\varphi_{T_i \prec T_j}$ commits to
ordering the actions of transaction $T_i$ before the conflicting actions of $T_j$.
If we observe $\{(b), (c), (d)\}$ as the query result, then an example determination is
\[
  I_{\mathsf{in}} \;\triangleq\;
  \bigl(
  \varphi_{T_1 \prec T_2}
  \wedge
  \varphi_{T_2 \prec T_3}
  \wedge
  \varphi_{T_3 \prec T_4}
  \wedge
  \varphi_{T_4 \prec T_Q}
  \wedge
  \varphi_{T_Q \prec T_5}
  \bigr)
\]
All predicates in $I_{\mathsf{in}}$ commute (each filters admissible serializations
by an independent pairwise constraint, so application order is irrelevant);
thus $I_{\mathsf{in}}$ forms a single
determination layer (depth~1).
The corresponding conditioned (depth--0) provenance is
\[
  \Prov_{Q \mid I_{\mathsf{in}}}(b)=x_1+x_2,\;
  \Prov_{Q \mid I_{\mathsf{in}}}(c)=x_3,\;
  \Prov_{Q \mid I_{\mathsf{in}}}(d)=x_4.
\]

Determination provenance pairs these expressions with the determination,
e.g., $\mathsf{DP}_Q(b) \triangleq (I_{\mathsf{in}},\, x_1+x_2)$.

In Appendix~\ref{sec:appendix-concurrency} we work this example more fully, including
the case where $d$ is absent from the query result.

\subsection{Determination Parsimony and Discriminating Power}
\label{sec:parsimony}

Recall that for a history $H$ and determination $I$, refinement induces a
quotient $O_H^I \big/ {\sim_{I,H}}$,
which records which admissible outcomes remain semantically distinguished after
conditioning (Section~\ref{sec:prelim}).

\paragraph{Discriminating power.}
Fix a specification $\Spec$.
A determination $I$ \emph{discriminates} between two outcomes
$o,o' \in \Spec(H)$ if $o \not\sim_{I,H} o'$.
If $o \sim_{I,H} o'$ for all resolving determinations,
then the specification lacks the expressive power to distinguish between the two outcomes.

\paragraph{Parsimonious commitment bases.}
Let $\mathcal{B}$ be a basis of commitments.
We call $\mathcal{B}$ a \emph{parsimonious} basis if:
(i) for every history $H$, some determination formed from $\mathcal{B}$
resolves $\Spec$ at $H$; and
(ii) for every $H$ and every pair of distinct admissible outcomes
$o \neq o' \in \Spec(H)$, there exists a resolving determination from
$\mathcal{B}$ that discriminates between them.

Intuitively, a parsimonious basis induces quotients that are no coarser than
necessary: any semantic distinction that exists in principle can be exposed by
some admissible determination.

\paragraph{Parsimony and robustness.}
Parsimony is a property of the \emph{commitment basis}, not of individual queries.
The commitment basis determines which semantic distinctions are expressible, and
therefore which admissible outcomes can be distinguished or collapsed.
In practice, systems often adopt non-parsimonious bases—for reasons of efficiency,
policy, or usability—thereby intentionally abstracting away fine-grained
differences among admissible outcomes.

In the transactional setting, for example, predicates that commit only to an
equivalent serial schedule are deliberately coarse.
Some tuples—such as $b$ and $c$ in Example~\ref{ex:example}—are present under every
determination consistent with that abstraction.
Determination provenance makes this explicit by showing which observations are
invariant across all determinations induced by the chosen commitment basis; these
observations require sharper ordering predicates to distinguish. But in many cases,
including transaction scheduling, such invariance is a side-effect of a semantically
natural commitment basis.

% This invariance is not accidental.
% Observations are robust precisely when the available commitment vocabulary
% never discriminates between the outcomes on which they would differ.
This perspective motivates questions studied in
Section~\ref{sec:robustness}, where this invariance (which we call \emph{robustness}) is defined
\emph{relative to a fixed commitment basis}, rather than as an absolute semantic
property.

\subsection{Summary}
Once a determination is fixed, provenance is classical: it is
computed over the conditioned specification $\Spec \mid I$.
The following sections instantiate this framework in two settings with very
different sources of semantic ambiguity: negation in Datalog and scheduling in
transactional systems.


% ================================================================
\section{Layer Certificates and Determination Provenance}
\label{sec:certificates}
% ================================================================

Fix a specification $\Spec$ and a determination
$I=\varphi_1\wedge\cdots\wedge\varphi_k$ with layering $L_1,\ldots,L_k$.

A \emph{provenance query} is a function $q:(\Spec\mid I)\to A$,
where $A$ is an answer domain (e.g., Boolean, explanations, dependency sets),
including classical why/why-not provenance and the robustness queries of
Section~\ref{sec:robustness}.
Such queries depend only on the restrictions that a determination places on the
space of admissible outcomes.
A \emph{certificate} summarizes these restrictions.

\begin{definition}[Certificate]
  Let $\varphi$ be a commitment.
  A certificate for $\varphi$, written $\varphi\models C$, is a predicate
  $C:O\to\{\mathsf{true},\mathsf{false}\}$ such that for all histories $H$ and
  nonempty $S\subseteq O$,
  \[
    \varphi(H,S)\subseteq\{\,o\in S\mid C(o)\,\}.
  \]
\end{definition}

Certificates are sound over-approximations: they may admit false positives, but
never exclude outcomes admitted by a determination.
Whether such lossiness is acceptable depends on the provenance queries being asked.
To formalize this dependence, write
\[
  \Spec_i \triangleq \Spec \mid (\varphi_1\wedge\cdots\wedge\varphi_i),
  \qquad
  J \triangleq \varphi_{i+1}\wedge\cdots\wedge\varphi_k .
\]

\begin{definition}[Sufficient Cert.]
  Let $\chi_{C_i}$ denote the commitment induced by certificate $C_i$, defined by
  $\chi_{C_i}(H,S) \triangleq \{\,o \in S \mid C_i(o)\,\}$.
  A certificate $C_i$ is \emph{sufficient for $(J,\mathcal{Q})$} if for all histories $H$,
  \[
    (\Spec_i\mid J)(H)=((\Spec\mid\chi_{C_i})\mid J)(H),
  \]
  and for all $q\in\mathcal{Q}$,
  \[
    q((\Spec_i\mid J)(H)) = q(((\Spec\mid\chi_{C_i})\mid J)(H)).
  \]
\end{definition}

\noindent
Thus, certificates may be lossy, but only insofar as they preserve both further
semantic resolution and the answers to the intended provenance queries.

\begin{definition}[Layer Cert.]
  A \emph{layer certificate} is a certificate $C_i$ for layer $L_i$ that is sufficient
  for $(J,\mathcal{Q})$.
\end{definition}

\begin{theorem}[Query Monotonicity]
  If $\mathcal{Q}_1\subseteq\mathcal{Q}_2$, then every certificate sufficient for
  $(J,\mathcal{Q}_2)$ is sufficient for $(J,\mathcal{Q}_1)$.
\end{theorem}
\begin{proof}[Proof sketch]
  Sufficiency requires that $C_i$ preserve both further semantic resolution and
  the answers to all queries in the target family.
  If $C_i$ is sufficient for $(J,\mathcal{Q}_2)$, then for every $q \in \mathcal{Q}_2$
  and every history $H$,
  $q((\Spec_i \mid J)(H)) = q(((\Spec \mid \chi_{C_i}) \mid J)(H))$.
  Since $\mathcal{Q}_1 \subseteq \mathcal{Q}_2$, this equality holds \emph{a fortiori}
  for every $q \in \mathcal{Q}_1$.
  The resolution-preservation condition is independent of $\mathcal{Q}$ and
  therefore inherited directly.
\end{proof}

\noindent
Thus, expanding the space of provenance questions can only require retaining
additional semantic information, never less.
Appendix~\ref{sec:appendix-datalog} illustrates this effect for semantic-change
queries under negation; Appendix~\ref{sec:appendix-concurrency} does so for
robustness queries under transactional ambiguity.


\paragraph{Summary.}
Layer certificates provide a principled interface between semantic resolution and
provenance, making explicit which semantic distinctions must be retained for a
given family of provenance queries.
Concretely: in the Datalog sandbox, a sealed set of atoms suffices for
derivational provenance but not for semantic-change queries
(Appendix~\ref{sec:appendix-datalog}); in the transactional sandbox, ordering
certificates suffice under serializability but snapshot structure must be
preserved for SI-to-SER preservation queries
(Appendix~\ref{sec:appendix-concurrency}).
This reframes a familiar tradeoff in systems lineage: tracing systems such as
Dapper retain full histories to support arbitrary diagnostics~\cite{sigelman2010dapper},
corresponding to certificates that preserve all future queries.
Determination provenance exposes this choice explicitly and enables more
parsimonious lineage when the intended query family is known
(Appendix~\ref{app:systems-lineage}).


\begin{figure*}[t]
  \centering
  \footnotesize
  \setlength{\tabcolsep}{6pt}
  \renewcommand{\arraystretch}{1.15}

  \begin{tabular}{@{}p{0.12\linewidth} p{0.65\linewidth} p{0.12\linewidth}@{}}
    \toprule
    \multicolumn{3}{c}{\textbf{(A) Datalog negation semantics}}           \\
    \midrule
    \textbf{Case} & \textbf{Resolving determination(s)} & \textbf{Depth} \\
    \midrule
    Stratified    &
    $\varphi_{\seal}(S_1) \seq \varphi_{\seal}(S_2) \seq \cdots \seq \varphi_{\seal}(S_k)$
                  & $k$                                                   \\

    Well-founded  &
    $\varphi_{\seal}(S_1) \seq \cdots \seq \varphi_{\seal}(S_k)
      \seq \varphi_{\seal}(U_1) \seq \varphi_{\seal}(U_2) \seq \cdots$
                  & unbounded                                             \\

    Stable        &
    $(\varphi_{\seal}(S_1) \seq \cdots \seq \varphi_{\seal}(S_k) \seq I^{(r)})
      \;\text{or}\;
      (\varphi_{\seal}(S_1) \seq \cdots \seq \varphi_{\seal}(S_k) \seq I^{(s)})$
                  & $k{+}1$                                               \\
    \bottomrule
  \end{tabular}

  \vspace{1em}

  \begin{tabular}{@{}p{0.23\linewidth} p{0.54\linewidth} p{0.12\linewidth}@{}}
    \toprule
    \multicolumn{3}{c}{\textbf{(B) Transactions: serializability and snapshot isolation}} \\
    \midrule
    \textbf{Case}                & \textbf{One resolving determination} & \textbf{Depth} \\
    \midrule
    SER robust $b$               &
    $\{\varphi_{T_1 \prec T_Q},\varphi_{T_Q \prec T_2},\varphi_{T_2 \prec T_3},
      \varphi_{T_3 \prec T_4},\varphi_{T_4 \prec T_5}\}$
                                 & $1$                                                    \\
    %
    SER contingent $d$ (why)     &
    $\{\varphi_{T_1 \prec T_2},\varphi_{T_2 \prec T_3},\varphi_{T_3 \prec T_4},
      \varphi_{T_4 \prec T_Q},\varphi_{T_Q \prec T_5}\}$
                                 & $1$                                                    \\
    %
    SER contingent $d$ (why-not) &
    $\{\varphi_{T_1 \prec T_2},\varphi_{T_2 \prec T_3},\varphi_{T_3 \prec T_4},
      \varphi_{T_4 \prec T_5},\varphi_{T_5 \prec T_Q}\}$
                                 & $1$                                                    \\
    %
    SI (write-skew $e/f$)        &
    $\varphi_{\snap(T_1)=\emptyset} \seq \cdots \seq
      \varphi_{\snap(T_7)=\emptyset}
      \seq \varphi_{T_4 \in \snap(T_5)}
      \seq \varphi_{\snap(T_Q)=\{T_1,\dots,T_7\}}
      \seq \varphi_{\FCW}(\{T_1,\dots,T_7\})$
                                 & $O(|H|)$                                               \\
    \bottomrule
  \end{tabular}

  \caption{
    Determination summaries for examples.
    Depth counts non-commuting stages of semantic commitment; branching among
    incompatible determinations does not increase depth.
    (\textbf{A}) For Datalog with negation, different semantics induce different
    resolution structures: stratified and well-founded semantics admit a single
    resolving determination, while stable semantics admits multiple incompatible
    ones.
    (\textbf{B}) For serializable transactions, many serial determinations may
    justify the same outcome; we therefore display one resolving determination
    as a witness, since only existence (not enumeration) matters.
    The SI row abbreviates snapshot choices ($\snap(\cdot)$) and the
    first-committer-wins rule ($\FCW$).
    Details appear in
    Appendices~\ref{sec:appendix-datalog} and~\ref{sec:appendix-concurrency}.
  }
  \label{fig:det-tables}
\end{figure*}

% ================================================================
\section{Sandbox I: Datalog with Negation as Semantic Ambiguity}
\label{sec:negation}
% ================================================================

We use Datalog with negation as a controlled setting in which to study
\emph{semantic ambiguity} within a fixed specification.
Our goal is not to propose new semantics, but to expose how classical negation
semantics differ in the \emph{structure} of the semantic commitments they require,
and how those commitments are reflected in determination provenance.

Even for a fixed program and extensional input, negation can delay semantic
commitment.
Whether an atom holds may depend on the absence of other atoms whose status is
itself unresolved.
Determination provenance makes this explicit by separating ordinary derivational
consequences—true in all admissible outcomes—from \emph{commitments}
that rule out some outcomes and thereby resolve ambiguity.

\begin{example}[Datalog with negation]
  \label{ex:datalog-negation}
  Consider the following program $P$ over EDB facts in a history $H$.
  For notational economy we use a fixed constant $c$ and take all predicates unary, so the EDB fact is $a(c)$ and the derived atoms are $p(c),q(c),r(c),s(c)$.
  \begin{align*}
    p(c) & \leftarrow a(c).      & r(c) & \leftarrow \neg s(c). \\[1mm]
    q(c) & \leftarrow \neg p(c). & s(c) & \leftarrow \neg r(c).
  \end{align*}
  Assume $H$ contains the single extensional fact $a(c)$.
  The pair $(P,H)$ contains both a stratified fragment
  ($q(c)$ depends negatively on $p(c)$, and $p(c)$ depends positively on $a(c)$)
  and a cycle through negation ($r \leftrightarrow \neg s$).
  We use this single instance to exhibit, side-by-side, three distinct structures
  of determination provenance.
\end{example}

We examine three standard semantics for negation—stratified, well-founded, and
stable—and record the corresponding resolving determinations.
Figure~\ref{fig:det-tables}(\textbf{A}) summarizes these determinations and
their \emph{determination depth}, defined as the length of a non-commuting chain
of semantic commitments.
When a semantics admits multiple incompatible resolutions, we list them;
when it admits a unique resolution, we record that determination alone.
Here $S_i$ denotes the set of atoms whose truth is sealed at stratum~$i$, and
$U_i$ denotes the unfounded set eliminated at alternation~$i$ in the
well-founded construction.
Full derivations and conditioned provenance appear in
Appendix~\ref{sec:appendix-datalog}.

\subsection{Negation Semantics as Determination Shapes}

We illustrate these determination structures using the single running program of
Example~\ref{ex:datalog-negation}, which contains both a stratified fragment
($p,q$) and a cycle through negation ($r \leftrightarrow \neg s$).
Depth counts only non-commuting stages of semantic commitment; branching among
incompatible determinations does not increase depth.

\paragraph{Stratified semantics.}
In stratified programs, negation flows strictly from higher to lower strata.
Resolving ambiguity therefore requires committing that each stratum has
\emph{stabilized}: no additional facts for predicates in that stratum will appear.
In the example, the lower stratum defining $p(c)$ must be sealed before $q(c)$ can
be interpreted negatively.

These sealing predicates do not commute: stabilization of stratum~$i$ is a
prerequisite for applying negation commitments in stratum~$i{+}1$.
Accordingly, a program with $k$ strata has determination depth exactly~$k$,
corresponding to the layered determination
\[
  \varphi_{\seal}(S_1) \seq \varphi_{\seal}(S_2) \seq \cdots \seq \varphi_{\seal}(S_k)
\]
% [added \seq]
Figure~\ref{fig:det-tables}(\textbf{A}, Stratified) schematically represents
this chain; the explicit predicates appear in
Appendix~\ref{sec:appendix-datalog}.

\begin{proposition}[Stratified determination depth]
  \label{prop:stratified-depth}
  A stratified Datalog program with $k$ strata has determination depth
  exactly~$k$ under the sealing commitment basis.
\end{proposition}
\begin{proof}[Proof sketch]
  Each stratum $S_i$ requires a sealing commitment $\varphi_{\seal}(S_i)$ that
  commits to the completeness of atoms in $S_i$.
  Sealing $S_i$ is a prerequisite for evaluating negation in $S_{i+1}$, so
  $\varphi_{\seal}(S_i)$ and $\varphi_{\seal}(S_{i+1})$ do not commute.
  Hence $k$ layers are necessary.
  Within each stratum, all sealing commitments commute (they act on disjoint
  predicate sets), so $k$ layers are also sufficient.
\end{proof}

\paragraph{Well-founded semantics.}
Well-founded semantics applies to programs with negation cycles, such as the
$r/s$ fragment of the example.
Rather than choosing among total models, it alternates between positive
derivation and elimination of unfounded sets, producing a three-valued outcome
in which atoms may be true, false, or undefined.

In determination terms, this corresponds to an \emph{alternating sequence} of non-commuting
sealing commitments layered after the stratified prefix:
\[
  \varphi_{\seal}(S_1) \seq \cdots \seq \varphi_{\seal}(S_k)
  \seq \varphi_{\seal}(U_1) \seq \varphi_{\seal}(U_2) \seq \cdots
\]
% [added \seq]
The number of alternations is not bounded in general, so well-founded semantics
does not admit a fixed finite determination depth.
Figure~\ref{fig:det-tables}(\textbf{A}, WFS) records the limit determination
for the example instance.

\paragraph{Stable model semantics.}
Stable semantics insists on total two-valued models.
Stable-model truth assignments rely on stratification to give negation a
closed-world meaning and therefore do not commute with stratification
commitments.
After sealing the stratified prefix of the program, a global hypothesis can be
made about the truth of atoms in the negation cycle (e.g., $r$ or $s$ in the
example) and checked for self-consistency.

Each stable model corresponds to a distinct resolving determination, and these
determinations are mutually incompatible.
Within any one stable determination, however, the commitments are logically
simultaneous and commute.
Thus stable semantics has determination depth~$k{+}1$, corresponding to
\[
  \bigl(\varphi_{\seal}(S_1) \seq \cdots \seq \varphi_{\seal}(S_k)\bigr)
  \seq I^{(r)}
  \quad\text{or}\quad
  \bigl(\varphi_{\seal}(S_1) \seq \cdots \seq \varphi_{\seal}(S_k)\bigr)
  \seq I^{(s)}
\]
% [added \seq]
despite potentially exponential branching across alternatives
(Figure~\ref{fig:det-tables}(\textbf{A}, Stable)).

\paragraph{Takeaway.}
The three semantics differ not only in which outcomes they admit, but in the
\emph{structure} of semantic resolution.
Stratified semantics requires a linear chain of stabilization commitments;
well-founded semantics requires an alternating, potentially unbounded sequence;
stable semantics branches among shallow but incompatible determinations after
stratification.
When negation is acyclic, all three semantics coincide extensionally, but
determination provenance still records the sealing commitments required to
justify negation.
Full derivations appear in
Appendix~\ref{sec:appendix-datalog}.

% ================================================================
\section{Sandbox II: Transactions as Systems Ambiguity}
\label{sec:transactions}
% ================================================================

We now turn to transactional systems.
As in Sandbox~I, our goal is not to propose new correctness criteria, but to use a
well-studied setting to expose how \emph{semantic ambiguity} arises within a fixed
specification, and how different resolution regimes induce different
\emph{determination structures}.

In transactional systems, ambiguity arises not from negation but from
\emph{concurrency}.
Transactions race to read and write shared state, and the execution history does
not uniquely determine which effects are observed.
Determination provenance explains how query results depend not only on data, but
on the semantic commitments used to resolve that concurrency.

\subsection{Concurrency as Semantic Ambiguity}

A transactional workload consists of a finite set of transactions issuing reads
and writes and eventually committing.
The history records these operations and commit events, but—crucially—does not
determine a unique order or visibility relation among transactions.
Even with complete knowledge of all operations, multiple outcomes may be
admissible.

The key observation is that transaction order is not a primitive feature of the
execution; it is a \emph{semantic commitment} used to justify an outcome.
This mirrors the role of stabilization commitments in Sandbox~I: before committing
to an order or visibility decision, certain conclusions are simply not determined.
Classical data provenance implicitly assumes such commitments have already been
made; determination provenance makes them explicit.

\subsection{Running Example}

We reuse the running example from the introduction and
Appendix~\ref{sec:appendix-concurrency}, focusing on the tuple $d$ inserted by
$T_4$ and deleted by $T_5$.
The history does not determine whether $d$ is visible to the query transaction
$T_Q$.
That depends on how concurrency is resolved.

As in Sandbox~I, this single example suffices to exhibit multiple determination
shapes under different semantics.

\subsection{Serializability as Determination Shape}

Under serializability, admissible outcomes are those explainable by some total
order of transactions consistent with the history.
Each outcome corresponds to at least one such serialization.

\paragraph{Commitments.}
The primitive commitments in this setting are \emph{ordering
  predicates}.
For transactions $T_i$ and $T_j$, the predicate $\varphi_{T_i \prec T_j}$ commits
that $T_i$ precedes $T_j$ in the serialization.
A determination is a finite conjunction of such predicates, inducing a partial
order; a fully resolving determination induces a total order.

Crucially, ordering predicates commute.
Committing to $T_1 \prec T_2$ and to $T_2 \prec T_3$ can be done in either order and
has the same effect.
As a result, serializability has determination depth~1: all semantic commitments
belong to a single commutative stage, which we can write schematically as
\[
  \varphi_{T_1 \prec T_2}, \ldots, \varphi_{T_i \prec T_j}
\]
or, exposing the single layer,
$\{\varphi_{T_i \prec T_j}\}_{i,j}$,
% [added \seq]
with no further sequencing required.

\paragraph{Example.}
Under the determination
$I_{\mathsf{in}} \triangleq
  \varphi_{T_1 \prec T_2 \prec T_3 \prec T_4 \prec T_Q \prec T_5}$,
the insert precedes the query and the delete follows it.
Conditioning on $I_{\mathsf{in}}$ yields a resolved outcome in which $d$ is present,
with classical provenance $\Prov_{Q \mid I_{\mathsf{in}}}(d)=x_4$.

Under the alternative determination
$I_{\mathsf{out}} \triangleq
  \varphi_{T_1 \prec T_2 \prec T_3 \prec T_4 \prec T_5 \prec T_Q}$,
the delete precedes the query and $d$ is absent.
Prior to fixing a determination, there is no single classical provenance question
for $d$.
Determination provenance records the semantic commitments under which each
explanation applies.

Figure~\ref{fig:det-tables}(\textbf{B}, SER) summarizes representative resolving
determinations.

\paragraph{Same observation, different derivational provenance.}
Even when multiple serializations yield the same observable outcome, they need not
agree on derivational provenance.
For instance, adding a transaction
$T_{-2}:\ \textsf{delete } S(2,b)$
yields serializations in which $b$ is observed with provenance $x_1$ or $x_1+x_2$,
depending on whether $T_{-2}$ precedes or follows $T_Q$.
Determination provenance distinguishes these semantic explanations even when the
query result is identical.

\subsection{Snapshot Isolation as Determination Shape}

Snapshot isolation (SI) resolves concurrency using a different basis of semantic
commitments.
Rather than committing to a total order, SI assigns each transaction a snapshot of
the database state it reads from and enforces a global
first-committer-wins constraint on writes.

\paragraph{Commitments.}
An SI determination consists of:
(i) snapshot predicates specifying, for each transaction, which committed effects
it observes; and
(ii) a global write-conflict admissibility predicate.
These commitments do not collapse to a total order and must be made separately for
potentially many concurrent transactions.

Unlike serializability, these predicates do not commute in general: the snapshot
visible to one transaction may depend on which earlier transactions have already
been assigned snapshots, because the first-committer-wins rule can render a
snapshot choice infeasible once a conflicting assignment is fixed.
In the worst case, each of the $n$ transactions in the history requires an
independent, non-commuting snapshot commitment, yielding determination
depth $O(|H|)$.
Schematically:
\[
  \varphi_{\snap(T_{i_1})} \seq \varphi_{\snap(T_{i_2})} \seq \cdots \seq \varphi_{\FCW}
\]
% [added \seq]
where later snapshot choices depend on earlier ones.

Figure~\ref{fig:det-tables}(\textbf{B}, SI) records one such resolving
determination for the running example.
Once an SI determination is fixed, the specification again becomes single-valued,
and provenance reduces to ordinary depth--0 derivational provenance.
What differs from serializability is the \emph{structure and depth} of the semantic
resolution required to reach that point.

\paragraph{Takeaway.}
Both serializability and snapshot isolation admit multiple admissible outcomes for
the same transactional history.
What distinguishes them is not the presence of ambiguity, but the
\emph{structure of semantic resolution}.
Serializability resolves ambiguity via a single commutative layer of ordering
commitments.
Snapshot isolation requires a richer, deeper structure of non-commuting snapshot
and conflict commitments.
Determination provenance makes these differences explicit, while classical
provenance applies only after resolution is complete.

% ================================================================
\section{Change and Robustness Across Determinations}
\label{sec:robustness}
% ================================================================

Determination provenance records not only how observations are derived, but the
semantic commitments under which those derivations are valid.
This section formalizes two such families of questions and grounds them in the
examples developed earlier.
Here, \emph{observations} refer to observable query answers or derived facts—i.e., elements of the outcome set $O$—as defined in Section~\ref{sec:prelim}.

% \subsection{Intra- vs.\ Inter-Specification Questions}

We distinguish two classes of provenance questions.
\emph{Intra-specification} questions fix a specification $\Spec$ and range over its
resolving determinations, asking which observations hold under different
admissible resolutions of the \emph{same} specification.
\emph{Inter-specification} questions compare observations across distinct
specifications $\Spec$ and $\Spec'$, asking which observations persist when the
specification itself changes—for example, when a transactional system moves from
serializability to snapshot isolation, or when a logic program adopts a different
negation semantics.
Both question families are expressed using determination provenance; they simply
quantify over different spaces.

\subsection{Intra-Specification Queries: Robustness}

Robustness is an intra-specification property.
Fix a specification $\Spec$ and a class $\mathcal{I}$ of resolving determinations.
An observation $o$ is \emph{robust} with respect to $\mathcal{I}$ iff
\[
  \forall I \in \mathcal{I}.\; o \in (\Spec \mid I)(H),
\]
where (recall Section~\ref{sec:prelim}) each $I \in \mathcal{I}$ is resolving, i.e., $\Spec \mid I$ is single-valued for each history.

Equivalently, $o$ does not depend on semantic commitments that vary across
$\mathcal{I}$.
In the lens of Section~\ref{sec:parsimony}, robustness is evidence of a lack of
parsimony in the commitment basis.

\begin{observation}[Robustness as non-discrimination]
  \label{obs:robustness-parsimony}
  An observation $o$ is robust with respect to $\mathcal{I}$ if and only if
  no commitment in the basis discriminates between admissible outcomes that
  agree on $o$.
  Formally, $o$ is robust iff for every $I, I' \in \mathcal{I}$,
  the resolved outcomes $(\Spec \mid I)(H)$ and $(\Spec \mid I')(H)$ agree on
  whether $o$ holds.
  Robustness therefore characterizes exactly the observations that lie in the
  intersection of all resolved outcomes induced by $\mathcal{I}$.
\end{observation}

\paragraph{Example: robustness under serializability.}
In the transactional example of
Appendix~\ref{sec:appendix-concurrency}, the history admits multiple serial orders.
Commitments record only ordering commitments between transactions.
Under all such orders, the query result includes the values $b$ and $c$, while the
presence of $d$ depends on whether the insert $T_4$ is ordered before the query
and the delete $T_5$ after it.

Our commitment basis captures the level of discrimination natural to
serializability: it distinguishes executions by transaction order, but not by
finer-grained, value-specific effects.
This exposes that $b$ and $c$ are robust under serializability, while $d$ is not.
The lack of parsimony in the commitment basis here is both meaningful (it exposes artifacts of serializability)
and practical (it avoids an explosion of ordering predicates).

\subsection{Inter-Specification Queries: Preservation under Semantic Change}
\label{sec:semantic-change}

Semantic change is an inter-specification question: it compares which observations
are admitted under different specifications for the same history.

\paragraph{Preservation.}
Given specifications $\Spec$ and $\Spec'$, outcome $o$, and
sets $\mathcal{I}$ and $\mathcal{I}'$ of all resolving determinations
for $\Spec$ and $\Spec'$ respectively,
$o$ is \emph{preserved} from $\Spec$ to $\Spec'$ iff
\[
  \forall I \in \mathcal{I},\;
  \bigl(o \in (\Spec \mid I)(H)\bigr)
  \;\Rightarrow\;
  \exists I' \in \mathcal{I}'.\;
  \bigl(o \in (\Spec' \mid I')(H)\bigr).
\]

\paragraph{Serializable outcomes under snapshot isolation.}
Snapshot isolation admits a larger space of determinations than serializability.
Preservation therefore reduces to the following question:
\emph{given an SI determination, does there exist a serializable determination
  that admits the same outcome?}

Determination provenance provides concrete syntax for this question.
An SI determination records snapshot choices and commit-admissibility constraints.
From these predicates, one can check whether the induced visibility and conflict
relations admit a topological sort constituting a serial execution.
If so, the outcome is preserved when strengthening from SI to serializability;
otherwise it is inherently non-serializable.
Appendix~\ref{sec:appendix-concurrency} works through a detailed example.

This reframes preservation as a \emph{semantic validation problem} over
determinations rather than as a comparison of executions.
Mixed-consistency execution policies—e.g., allowing SI but ignoring
any non-serializable outcome rows—can therefore be expressed and enforced directly in
terms of determination provenance, without re-executing the workload.

\paragraph{Summary.}
Robustness identifies conclusions invariant under the semantic distinctions a
system chooses to make.
Semantic change compares conclusions across different semantic regimes.
Determination provenance provides a common framework for expressing both, while
making explicit the limits of discrimination imposed by the chosen
commitments.
When the space of determinations is finite or otherwise measurable, these
qualitative questions admit quantitative variants—for example, what fraction of
admissible determinations justify a conclusion.

% ================================================================
\section{Related Work}
\label{sec:related}
% ================================================================

\noindent\textbf{Data provenance and explanations.}
Algebraic provenance provides a uniform framework for explaining query
conclusions under fixed, resolved semantics by annotating derivations with
elements of a commutative semiring
\cite{green2007provenance,suciu2011probabilistic,dalvi2012probabilistic,olteanu2015factorized,senellart2018provenance}.
A related body of work studies explanations and causality for query answers,
including why-not provenance, causal responsibility, and counterfactual
explanations
\cite{buneman2001why,halpern2001causes,meliou2010causality,glavic2015survey}.

From the perspective of this paper, these approaches operate at
determination depth~0: semantic ambiguity has already been resolved, and
provenance explains derivations relative to a fixed semantics.
This includes work on incremental view maintenance and \emph{homeostasis},
which identifies queries whose answers are invariant under admissible input
changes and therefore require no additional semantic commitments to maintain
\cite{gehrke2002homeostasis}.
In our framework, homeostasis corresponds precisely to robustness at
determination depth~0: the query answer is invariant across all admissible
outcomes of the specification, so no commitment is needed to justify it.
Determination provenance generalizes this by explaining which commitments are
required when invariance fails, and how classical provenance depends on those
commitments once resolution has occurred.

\medskip
\noindent\textbf{Negation and non-monotonic semantics.}
Classical semantics for Datalog with negation provide well-understood mechanisms
for resolving semantic ambiguity, with established relationships under syntactic
restrictions.
We do not propose new semantics.
Instead, Section~\ref{sec:negation} uses Datalog with negation as a controlled
setting in which resolution strategies are explicit and formal, making it
possible to study how different semantic commitments induce distinct
determination structures and provenance behavior.

\medskip
\noindent\textbf{Transactional and systems provenance.}
Provenance has been studied in transactional, workflow, and database
settings under an assumed execution or isolation level
\cite{bernstein1987concurrency,adya1999weak,cerone2015consistency,cheney2009provenance,freire2008provenance}.
In these settings, semantic resolution is implicit: provenance explains results
relative to a chosen execution, snapshot, or schedule.
From our perspective, such approaches again operate at determination depth~0.

A distinct body of systems work studies execution provenance or systems lineage,
recording causal dependencies among processes, messages, and events to explain
observed behavior
\cite{seltzer2005provenance,bates2015trustworthy,sigelman2010dapper}.
These approaches take a maximalist view of provenance, capturing entire execution
histories to account for ambiguity in scheduling and message interleaving.
Determination provenance provides a semantic foundation for structuring such
systems provenance: rather than committing to a single reconstructed execution,
it makes explicit the semantic commitments required for a diagnostic conclusion
to hold across admissible executions, supporting principled notions of parsimony
and robustness.
Recent systems work on mixed isolation regimes, such as Milano et al.’s MixT
\cite{milano2020mixt}, addresses related concerns operationally; our contribution
is a semantic and explanatory account that makes such mixtures explicit and
checkable.

\medskip
\noindent\textbf{Causal analysis of logs and executions.}
Recent work on causal analysis of system logs studies how diagnostic conclusions
depend on modeling choices made when constructing causal explanations from
execution traces, including discretization, causal graph construction,
sensitivity to unobserved variables, and summarization
\cite{markakis2024sawmill,chen2025causalens,markakis2024press,zeng2024validity,lai2025seercuts,levy2025causumx}.
These systems typically treat logs as causal histories and focus on summarizing,
querying, or stress-testing inferred causal structure.

From the perspective of this paper, such modeling choices can be understood as
semantic commitments that restrict the space of admissible explanations.
Determination provenance does not perform causal inference; rather, it provides
a framework for making these commitments explicit and for reasoning about how
conclusions depend on them.
In particular, the robustness and preservation questions studied in
Section~\ref{sec:robustness} suggest a natural connection: diagnostic conclusions
may be robust across a class of causal determinations, or may depend critically
on specific causal assumptions.
Clarifying this connection—by treating causal graph construction and
summarization as forms of semantic resolution over ambiguous histories—is a
promising direction for future work.

\medskip
\noindent\textbf{Determination complexity.}
This work builds on a prior complexity theory of determination
introduced in~\cite{hellerstein2026interpretation}, which characterizes the cost,
depth, and coordination requirements of semantic resolution.
The present paper focuses instead on explanation.
Determinations are assumed given, and determination provenance explains how
observations depend on them.
Connections to descriptive complexity, fixpoint alternation, oracle models, and
coordination are discussed in detail in~\cite{hellerstein2026interpretation}.

% ================================================================
\section{Conclusion}
\label{sec:conclusion}
% ================================================================

This paper reframes provenance around a simple observation: many modern systems do
not produce a single semantic outcome from a history, but a space of admissible
ones.
In such settings, explaining an observation requires accounting for how ambiguity
was resolved—not merely how derivations proceeded after the fact.
Determination provenance makes those semantic commitments explicit and
composable, recovering classical algebraic provenance once resolution is fixed
while retaining visibility into the assumptions that made resolution possible.

Our development shows that this perspective is not domain-specific.
In transactional systems, determination provenance explains how query
observations depend on ordering and isolation commitments.
In logic programming with negation, it exposes the semantic dependencies that
govern resolution.
Across domains, certificates provide a minimal interface for carrying forward
exactly the information required to support further reasoning.

Determination provenance also opens a broader research agenda.

\paragraph{Robustness and semantic sensitivity.}
Once observations are annotated with the determinations under which they hold, it
becomes natural to ask which are robust across determinations and which are
fragile.
This enables provenance questions about semantic stability, policy change, and
compatibility across regimes that lie beyond traditional why- and how-provenance.
Developing principled robustness measures—qualitative or quantitative—remains an
open direction.

\paragraph{Engineering determination provenance.}
While determinations are minimal by construction, the choice of
\emph{commitment basis} is not.
Designing effective commitment bases—balancing parsimony against the desire to
expose semantic sensitivity—raises new engineering questions at the boundary of
specification and provenance.
Coarser bases yield more robust observations but collapse distinctions; finer
bases increase discriminating power but can amplify fragility and cost, as
illustrated by isolation-levels in transactional systems (Section~\ref{sec:robustness}).
Understanding how to design specifications and commitment bases that make the
right semantic distinctions explicit, while supporting compact certificates and
efficient maintenance, is a central direction for future work.

\paragraph{Beyond databases.}
Semantic ambiguity and resolution arise well beyond data management.
Consistency models, coordination protocols, access-control systems, and machine
learning pipelines all exhibit underspecification followed by structured
commitment.
Determination provenance provides a vocabulary for reasoning across these
settings, suggesting a unifying foundation for provenance in ambiguous
computational systems.

The success of algebraic provenance lay in identifying the right semantic object
to explain.
Determination provenance extends that insight: when meaning itself depends on
resolution, provenance must explain the assumptions that made meaning possible.

\appendix
% ================================================================
\section{Appendix: Algebraic Structure After Semantic Resolution}
\label{app:algebra}
% ================================================================

This appendix clarifies the algebraic structure induced by determinations and
explains why classical algebraic provenance applies once semantic ambiguity has
been resolved.
The goal is not to develop a new provenance algebra, but to situate determination
provenance precisely relative to existing algebraic frameworks.

\subsection{Sequential Conditioning Versus Algebraic Combination}

Determinations refine a specification by applying commitments via
\emph{sequential conditioning}.
Given a specification $\Spec$ and commitments $\varphi_1,\varphi_2$, refinement
\[
  (\Spec \mid \varphi_1)\mid \varphi_2
\]
need not equal
\[
  (\Spec \mid \varphi_2)\mid \varphi_1,
\]
because predicates may be ambiguity-sensitive.
Accordingly, predicate composition across layers is generally non-commutative and
does not admit a semiring structure.

This distinction is essential.
Sequential conditioning eliminates admissible outcomes; it is a semantic operation,
not an algebraic one.
Commitments therefore do \emph{not} combine explanations, but restrict
the space in which explanations are defined.

\subsection{Commutative Structure Within a Layer}

Within a single determination layer, commitments commute by construction.
Fix a specification prefix $\Spec^{\ell}$ obtained after resolving all prior layers.
Let $L$ be a set of commitments applicable at that layer.

For $\varphi,\psi \in L$,
\[
  (\Spec^{\ell} \mid \varphi)\mid \psi
  \;=\;
  (\Spec^{\ell} \mid \psi)\mid \varphi .
\]
Thus commitments in $L$ form a \emph{commutative monoid} under conjunction, with the
identity given by the trivial commitment that performs no restriction.

This monoid is \emph{conditional}: the available commitments and their effect depend
on the determination prefix.
No global algebra over commitments is assumed or required.

\subsection{Special Case: Semiring Structure}

In special cases, the commutative monoid within a layer extends to a semiring.
This occurs when:
(i) commitments are pointwise over admissible outcomes; and
(ii) conjunction distributes over disjunction of outcomes.

These conditions coincide with those assumed in classical algebraic provenance.
Although many of our examples satisfy them, we do not exploit this structure in the
main text.
Our focus is on semantic resolution, not on algebraic optimization or factorization.

\subsection{Resolved Specifications and Monotone Evaluation}

The key transition occurs once a resolving determination $I$ has been fixed.

\begin{definition}[Resolved specification]
  A specification $\Spec \mid I$ is \emph{resolved} if for every history $H$,
  $(\Spec \mid I)(H)$ is a singleton.
\end{definition}

Resolution has two crucial consequences:

\paragraph{Single-valuedness.}
There are no remaining admissible alternatives.
Commitments no longer apply, and semantic conditioning ceases.

\paragraph{Monotonicity.}
All remaining variation is derivational.
Facts accumulate, explanations only grow, and no semantic commitments rule out
previously admissible outcomes.

Together, these properties place $\Spec \mid I$ in the standard setting of
monotone query evaluation.

\subsection{Why Semiring Provenance Applies at Layer~0}

Classical semiring provenance applies precisely to monotone evaluation over a fixed
instance.
Once determination is complete, this is exactly the situation we obtain.

\begin{proposition}[Resolved specifications admit semiring provenance]
  \label{prop:semiring}
  If $\Spec \mid I$ is resolved and the commitment basis is persistent
  (Definition~\ref{def:persistent}), then derivational provenance over
  $\Spec \mid I$ is captured by a commutative semiring.
\end{proposition}
\begin{proof}[Proof sketch]
  By Theorem~\ref{thm:resolution-implies-monotone}, $\Spec \mid I$ is
  history-monotone.
  Because it is also single-valued (resolved), $(\Spec \mid I)(H)$ is a singleton
  for every history $H$; no further commitments can eliminate outcomes, so semantic
  conditioning has ceased.
  Because it is monotone, the unique outcome grows monotonically with
  the history: $H_1 \sqsubseteq H_2$ implies that every fact derivable at $H_1$
  remains derivable at $H_2$.
  Under these two conditions, derivations compose exactly as in monotone query
  evaluation over a fixed instance: alternative derivations of the same fact
  combine by semiring addition, and joint use of multiple facts combines by
  semiring multiplication.
  This is precisely the algebraic regime established
  by Green et al.~\cite{green2007provenance} for positive relational queries.
  Although $\Spec \mid I$ need not syntactically be a Datalog program, it is
  algebraically equivalent, for provenance purposes, to monotone evaluation over
  a fixed instance, which suffices.
\end{proof}

\subsection{Relationship to Datalog and Prior Work}

Many of our examples are naturally expressed in Datalog, and classical provenance
results apply directly in those cases.
However, our claim is algebraic rather than syntactic:
determination provenance recovers the algebraic conditions required for semiring
provenance by explicitly discharging semantic ambiguity.

Seen this way, determination theory serves as a semantic prelude to algebraic
provenance.
It explains \emph{when} and \emph{why} classical provenance applies, rather than
assuming a fixed monotone semantics from the outset.

This perspective complements existing work on semiring provenance and factorized
representations by separating semantic resolution from algebraic explanation.


% \section{Appendix: Worked Examples}
% \label{sec:appendix-examples}

\section{Datalog$^{\neg}$ Worked Determination Provenance}
\label{sec:appendix-datalog}

This appendix instantiates determination provenance for the Datalog-with-negation
sandbox used throughout Section~\ref{sec:negation}.
We work through the same running example under stratified, well-founded, and stable
semantics, making explicit the commitments, their dependency structure,
and the resulting provenance annotations.
The purpose is to justify the determination shapes and depths summarized in
Figure~\ref{fig:det-tables}(\textbf{A}).

Throughout, we use \emph{sealing predicates}
$\varphi_{\seal}(X)$, which commit that the set of atoms $X$ is complete:
no additional atoms in $X$ will be derived in any admissible outcome.
Such predicates rule out outcomes and therefore constitute genuine semantic
commitments.

\subsection{Stratified Semantics}

\paragraph{Determination structure.}
In stratified semantics, predicates are partitioned into strata
$S_1,\dots,S_k$ such that negation flows strictly from higher to lower strata.
Resolving ambiguity requires committing that each stratum has stabilized before
negation may be applied in the next stratum.

For the running example (Example~\ref{ex:datalog-negation}), the stratified fragment
consists of the predicates defining $p(c)$ and $q(c)$.
Resolution therefore requires sealing the lower stratum containing $p(c)$.

\paragraph{Resolving determination.}
A resolving determination for the stratified fragment is
\[
  I_{\mathrm{strat}}
  \;\triangleq\;
  \varphi_{\seal}(S_1)
  \seq
  \varphi_{\seal}(S_2)
  \seq \cdots \seq
  \varphi_{\seal}(S_k),
\]
% [added \seq]
where each $S_i$ is the set of atoms whose truth is fixed at stratum~$i$.
These sealing predicates do \emph{not} commute: sealing $S_i$ is a prerequisite
for applying negation commitments in $S_{i+1}$.
Accordingly, the determination depth is exactly $k$, as recorded in
Figure~\ref{fig:det-tables}(\textbf{A}, Stratified).

\paragraph{Provenance.}
Once $I_{\mathrm{strat}}$ is fixed, the specification is fully resolved.
Derivational provenance applies without further semantic refinement.
For example,
\[
  \mathsf{DP}(p(c)) \;\triangleq\; (I_{\mathrm{strat}},\, \Prov(a(c))).
\]
Absences such as $\neg q(c)$ are derivational consequences of the sealed lower stratum
and require no additional commitments.

\subsection{Well-Founded Semantics}

\paragraph{Determination structure.}
Well-founded semantics (WFS) applies to programs with negation cycles.
It resolves ambiguity by alternating between positive derivation and the elimination
of \emph{unfounded sets}, yielding a three-valued outcome domain
$\{\mathbf{t},\mathbf{f},\mathbf{u}\}$.

Crucially, WFS still relies on the same stratified-prefix commitments as stratified
semantics: the lower strata must be sealed before negation is interpreted.
Beyond this prefix, resolution proceeds through an \emph{alternating sequence} of
additional sealing commitments that classify atoms as determined or undefined.

\paragraph{Resolving determination.}
A resolving determination under WFS can be written schematically as
\[
  I_{\mathrm{wfs}}
  \;\triangleq\;
  \varphi_{\seal}(S_1)
  \seq \cdots \seq
  \varphi_{\seal}(S_k)
  \seq
  \varphi_{\seal}(U_1)
  \seq
  \varphi_{\seal}(U_2)
  \seq \cdots,
\]
% [added \seq]
where each $U_i$ is an unfounded set eliminated at alternation~$i$.
These commitments are non-commuting and may require an unbounded number of
alternations.
Hence WFS has no fixed finite determination depth in general, as indicated in
Figure~\ref{fig:det-tables}(\textbf{A}, Well-founded).

For the running example, the limit of this process yields the unique partial
outcome
\[
  p(c)=\mathbf{t},\quad q(c)=\mathbf{f},\quad r(c)=\mathbf{u},\quad s(c)=\mathbf{u}.
\]

\paragraph{Provenance.}
Because WFS yields a single resolved outcome (in the lifted domain), there is no
branching among determinations.
Once $I_{\mathrm{wfs}}$ is fixed, provenance is again depth~0:
\[
  \mathsf{DP}(p(c)) \;\triangleq\; (I_{\mathrm{wfs}},\, \Prov(a(c))).
\]
The difference from stratified semantics lies not in multiplicity of outcomes, but
in the unbounded, alternating structure of the determination itself.

\subsection{Stable Model Semantics}

\paragraph{Determination structure.}
Stable semantics insists on total two-valued models.
Like WFS, it requires sealing the stratified prefix before negation is resolved.
Unlike WFS, it does not lift the outcome domain; instead it resolves cycles by
branching among incompatible global hypotheses.

\paragraph{Resolving determinations.}
For the running example, resolution begins with the same stratified-prefix sealing
\[
  \varphi_{\seal}(S_1)
  \seq \cdots \seq
  \varphi_{\seal}(S_k),
\]
% [added \seq]
After this prefix, the negation cycle admits two incompatible resolutions:
\[
  I^{(r)}
  \;\triangleq\;
  \bigl(
  \varphi_{\seal}(S_1)
  \seq \cdots \seq
  \varphi_{\seal}(S_k)
  \bigr)
  \seq
  (r(c)=\mathbf{t},\,s(c)=\mathbf{f}),
\]
\[
  I^{(s)}
  \;\triangleq\;
  \bigl(
  \varphi_{\seal}(S_1)
  \seq \cdots \seq
  \varphi_{\seal}(S_k)
  \bigr)
  \seq
  (s(c)=\mathbf{t},\,r(c)=\mathbf{f}).
\]
% [added \seq]

Within either determination, the remaining commitments are logically simultaneous
and commute.
Thus the determination depth is $k{+}1$: the $k$ stratification stages plus one
additional, non-commuting branching stage.
This is reflected in Figure~\ref{fig:det-tables}(\textbf{A}, Stable).

\paragraph{Provenance.}
Determination provenance must record which stable determination is assumed:
\[
  \mathsf{DP}(r(c)) \;\triangleq\; (I^{(r)},\,1),
  \qquad
  \mathsf{DP}(s(c)) \;\triangleq\; (I^{(s)},\,1).
\]
Only after this choice does ordinary derivational provenance apply.

\subsubsection{Semantic Change Across Negation Semantics}

Holding the program and extensional history fixed, different negation semantics
induce different spaces of resolving determinations.

For example, consider the program
\[
  p(c)\leftarrow a(c),\qquad
  q(c)\leftarrow \neg r(c),\qquad
  r(c)\leftarrow \neg q(c),
\]
with $H=\{a(c)\}$.
Under WFS, there is a unique resolving determination with
$q(c)=r(c)=\mathbf{u}$.
Under stable semantics, there are two incompatible determinations, one with
$q(c)=\mathbf{t}$ and one with $r(c)=\mathbf{t}$.

The conclusion $p(c)$ is preserved across semantics, while $q(c)$ and $r(c)$ are not.
Determination provenance makes this preservation (and non-preservation) explicit
and checkable.

\subsection{Certificates and Semantic-Change Queries}

Fix a specification $\Spec$ and a resolving determination $I$.
A \emph{provenance query} is a function
\[
  q : (\Spec \mid I) \to A,
\]
where $A$ is an answer domain (e.g., Boolean, explanations, dependency sets),
including classical why/why-not provenance and the semantic-change queries of
Section~\ref{sec:robustness}.
Such queries depend only on the restrictions that a determination imposes on
admissible outcomes, not on the full resolution process.

Commitments restrict admissible outcomes; a \emph{certificate}
summarizes these restrictions.
Formally, a certificate is a sound over-approximation: it may admit outcomes that
a full determination would exclude, but must exclude none that a determination
admits.
This lossiness is acceptable precisely because provenance queries are evaluated
only up to the distinctions they are designed to observe.

In the negation sandbox, a natural certificate for the stratified prefix is given
by a \emph{sealed set} of atoms.
Equivalently, this can be represented by a \emph{history antichain}: a frontier in
the partial order of derivations that fixes which facts are considered complete.
In a totally ordered execution, a single timestamp freezes the past; in a
partially ordered history, an antichain plays the same role, sealing exactly the
downward closure of a sub-history without imposing an artificial total order.

Under a fixed resolving determination, a minimal certificate sufficient for
derivational provenance of $p(c)$ need only preserve the sealing commitments for
the stratified prefix.
By contrast, answering a semantic-change query such as
``Would $q(c)$ become true under stable semantics?''
requires a strictly stronger certificate: it must preserve the unresolved
negation structure itself, not merely one resolution.

This illustrates the query-parametric nature of certificate sufficiency
(Section~\ref{sec:certificates}): supporting richer provenance questions strictly
increases the semantic information that must be retained.

% ===============================================================
\section{Transactional Worked Determination Provenance}
\label{sec:appendix-concurrency}
% ===============================================================

This appendix instantiates determination provenance for the transactional sandbox
used in Section~\ref{sec:transactions}.
As in the Datalog appendix, we make explicit the commitments,
their dependency structure, and the resulting provenance annotations.
The goal is to justify the determination shapes and depths summarized in
Figure~\ref{fig:det-tables}(\textbf{B}).

Throughout, commitments rule out admissible executions and therefore
constitute genuine semantic commitments.
Once a resolving determination is fixed, the specification becomes single-valued
and provenance reduces to ordinary derivational provenance.

% ===============================================================
\subsection{Serializability}
% ===============================================================

\paragraph{Determination structure.}
Under serializability, semantic ambiguity arises because a concurrent history does
not determine a unique transaction order.
Outcomes are justified by committing to a \emph{serialization} consistent with the
history’s conflicts.

The primitive commitments are \emph{ordering predicates}
$\varphi_{T_i \prec T_j}$, asserting that transaction $T_i$ precedes $T_j$ in the
serialization.
Conditioning on such a predicate eliminates all admissible outcomes whose
serializations violate it.

Crucially, ordering predicates commute: no ordering commitment enables or disables
another.
Semantic resolution therefore consists of a \emph{single stage} of commitment,
regardless of how many orderings are asserted.
Hence serializability has determination depth~1.

\paragraph{Data, query, and history.}
Let $S(k,y)$ be a relation and consider the query
$Q(y) \drule S(k,y)$.
The database is initially empty.
Five transactions execute:
\[
  \begin{array}{ll}
    T_1: \textsf{insert } S(1,b) \ (x_1) \qquad &
    T_2: \textsf{insert } S(2,b) \ (x_2)          \\
    T_3: \textsf{insert } S(3,c) \ (x_3) \qquad &
    T_4: \textsf{insert } S(4,d) \ (x_4)          \\
    T_5: \textsf{delete } S(4,d).               &
  \end{array}
\]
A query transaction $T_Q$ evaluates $Q$.

The history alone does not determine whether $d$ is visible to $T_Q$; this depends
on the relative ordering of $T_4$, $T_5$, and $T_Q$.

\paragraph{Resolving determinations.}
Two valid resolving determinations under serializability are:
\[
  I_{\mathsf{in}} \triangleq
  T_1 \prec T_2 \prec T_3 \prec T_4 \prec T_Q \prec T_5,
  \qquad
  I_{\mathsf{out}} \triangleq
  T_1 \prec T_2 \prec T_3 \prec T_4 \prec T_5 \prec T_Q.
\]
Both are conjunctions of ordering predicates and therefore arise in a single
semantic stage.

\paragraph{Provenance.}
Once a resolving determination is fixed, provenance is purely derivational.
Conditioned on $I_{\mathsf{in}}$, the resolved snapshot contains $b,c,d$:
\[
  \Prov_{Q \mid I_{\mathsf{in}}}(b)=x_1+x_2,\quad
  \Prov_{Q \mid I_{\mathsf{in}}}(c)=x_3,\quad
  \Prov_{Q \mid I_{\mathsf{in}}}(d)=x_4.
\]
Conditioned on $I_{\mathsf{out}}$, tuple $d$ is absent and classical why-not
provenance applies.

Determination provenance records the semantic assumption under which the
derivation holds; e.g.,
\[
  \mathsf{DP}_Q(d) \triangleq (I_{\mathsf{in}},\, x_4).
\]

\paragraph{Same observation, different derivational provenance.}
Even when multiple serializations yield the same observation, they need not agree
on derivational provenance.
Adding
$T_{-2}:\ \textsf{delete } S(2,b)$
yields executions in which $b$ is observed with provenance $x_1$ or $x_1+x_2$,
depending on the serialization.
Thus determination provenance distinguishes semantic explanations even when the
observable outcome is unchanged.

% ===============================================================
\subsection{Snapshot Isolation}
\label{sec:appendix-SI}
% ===============================================================

\paragraph{Determination structure.}
Snapshot isolation (SI) resolves concurrency using a different basis of semantic
commitments.
Rather than committing to a total order, an SI determination assigns each
transaction a \emph{snapshot} of committed transactions whose effects it reads,
together with a global write-conflict admissibility constraint.

In determination terms, this yields two kinds of predicates:
(i) snapshot predicates $\snap(T)=U$, committing that $T$ reads from the effects of
exactly the transactions in $U$, and
(ii) a global \emph{first-committer-wins} predicate $\varphi_{\FCW}$ that restricts
which snapshot assignments are admissible.

Unlike ordering predicates, snapshot commitments do not commute in general.
Each transaction may require an independent snapshot decision, yielding a staged
determination structure whose depth grows with the history.

\paragraph{Extended history.}
In addition to $T_1$–$T_5$, assume the database initially contains $S(5,e)$ and
$S(6,f)$, and the following transactions execute concurrently:
\[
  \begin{array}{ll}
    T_6: \textsf{read } S(5,e);\ \textsf{delete } S(6,f) \qquad &
    T_7: \textsf{read } S(6,f);\ \textsf{delete } S(5,e).
  \end{array}
\]
A query transaction $T_Q$ again evaluates $Q$.

\paragraph{Resolving determination.}
One admissible SI determination is:
\[
  \begin{aligned}
    I_{\mathrm{SI}} \triangleq {} &
    \snap(T_1)=\snap(T_2)=\snap(T_3)=\snap(T_4)=\snap(T_6)=\snap(T_7)=\emptyset \\
                                  & \seq\; \snap(T_5)=\{T_4\}                   \\
                                  & \seq\; \snap(T_Q)=\{T_1,\dots,T_7\}         \\
                                  & \seq\; \varphi_{\FCW}.
  \end{aligned}
\]
% [added \seq]

\paragraph{Outcome and provenance.}
Under $I_{\mathrm{SI}}$, the final state contains $b$ and $c$, but not $d,e,f$.
The query result is $\{b,c\}$ with derivational provenance:
\[
  \Prov_{Q \mid I_{\mathrm{SI}}}(b)=x_1+x_2,\qquad
  \Prov_{Q \mid I_{\mathrm{SI}}}(c)=x_3.
\]

\paragraph{Semantic change and serializability.}
Although $I_{\mathrm{SI}}$ is admissible under snapshot isolation, it does not admit
any serial execution.
The concurrent writes of $T_6$ and $T_7$ produce the classical write-skew anomaly:
a dependency cycle that cannot be linearized into a total order.

Determination provenance makes this explicit.
Restricting attention to SI determinations whose snapshot predicates admit a
topological sort recovers exactly the serializable outcomes.

\paragraph{Determination depth.}
Unlike serializability, which resolves ambiguity in a single commuting stage,
snapshot isolation may require an independent snapshot commitment for each
transaction.
Thus SI has worst-case determination depth $O(|H|)$, as recorded in
Figure~\ref{fig:det-tables}(\textbf{B}).
Once an SI determination is fixed, provenance again reduces to depth--0
derivational provenance.

% ===============================================================
\subsection{Certificates and Semantic-Change Queries}
% ===============================================================

As in the negation sandbox, provenance queries depend only on the restrictions that
a determination imposes on admissible outcomes.
A \emph{certificate} summarizes these restrictions and may safely over-approximate
the determination.

Under serializability, a certificate need only record a set of ordering predicates
sufficient to justify the observation.
Under snapshot isolation, however, answering semantic-change queries—such as
whether an SI outcome admits a serial explanation—requires preserving the snapshot
structure itself, not merely a particular resolved state.

As in Section~\ref{sec:appendix-datalog}, certificates are therefore
query-parametric.
Richer provenance questions require retaining strictly more semantic information,
reflecting the deeper and less-commutative structure of the underlying
determination.

% ================================================================
\section{Appendix: Determination Provenance for Systems Lineage}
\label{app:systems-lineage}
% ================================================================

This appendix illustrates determination provenance in a systems setting that
lies outside traditional database semantics.
The example is adapted from the distributed tracing model introduced by
Sigelman et al.~\cite{sigelman2010dapper}, abstracting its trace semantics rather
than its implementation details.

The example serves two purposes.
First, it shows that determination provenance extends the database tradition of
provenance to forms of \emph{systems lineage} that have historically resisted
formal treatment.
Second, it demonstrates how the parsimony of determinations---as minimal semantic
commitments---can refine systems lineage explanations by making explicit which
assumptions are required to justify diagnostic conclusions.

Unlike the previous appendices, this setting does not admit a fixed canonical
determination depth.
Semantic commitments may be staged or bundled depending on the diagnostic query,
but determination provenance still cleanly separates observable evidence from
the assumptions required to explain it.

\paragraph{Scope.}
The treatment below is deliberately illustrative rather than fully formal.
The Datalog and transactional sandboxes admit precise commitment bases because
their semantics are mathematically specified; distributed tracing, by contrast,
involves physical clocks and operational assumptions that resist clean
axiomatization.
We therefore describe commitments semi-formally, showing how the determination
provenance framework applies in principle while deferring a complete
formalization to future work on specific tracing systems.

\subsection{Trace Specification and Semantic Ambiguity}

A distributed trace consists of a set of \emph{spans}, each representing an
invocation of a service.
Each span $X$ has a client-side send and receive event
$(\mathsf{cs}_X,\mathsf{cr}_X)$ and a server-side receive and send event
$(\mathsf{sr}_X,\mathsf{ss}_X)$.
Client and server events occur on different hosts with unsynchronized clocks.

The trace specification $\Spec$ maps a history $H$ of timestamped span events to a
set of admissible executions $O$.
An execution assigns each span a concrete placement on a global timeline and
induces a partial order over spans.
The specification enforces only causal constraints:
\[
  \mathsf{cs}_X \prec \mathsf{sr}_X
  \qquad\text{and}\qquad
  \mathsf{ss}_X \prec \mathsf{cr}_X ,
\]
together with parent--child relationships recorded in the trace metadata.

Because clocks are unsynchronized, these constraints determine only \emph{bounds}
on server-side timestamps.
As a result, multiple executions---with different relative orderings and overlaps
of spans---are admissible under the same observed history.
Diagnostic conclusions such as critical-path analysis are therefore semantically
ambiguous prior to determination.

\subsection{Diagnostic Query: Critical Path}

We consider the diagnostic query:
\begin{quote}
  \emph{Which span lies on the critical path of the request?}
\end{quote}
Formally, the query maps an execution to the span (or set of spans) whose duration
determines the end-to-end latency of the root span.

Under a fixed execution, this query is deterministic.
Under the specification $\Spec$, however, different admissible executions may
yield different critical paths.
Without additional semantic commitments, the query is ill-posed.

\subsection{Commitments}

Determination resolves this ambiguity by committing to semantic assumptions that
fix a particular execution.
Relevant commitments include:
\begin{itemize}
  \item $\varphi_{\mathrm{align}(h_i,h_j)}$, committing to a concrete clock
        alignment between hosts $h_i$ and $h_j$ within admissible bounds;
  \item $\varphi_{\mathrm{order}(X,Y)}$, committing that span $X$ precedes span
        $Y$ when their relative order is otherwise unconstrained;
  \item $\varphi_{\mathrm{cp}=X}$, committing that $X$ is the critical-path
        witness \emph{given the preceding alignment and ordering commitments}.
\end{itemize}

These predicates are ambiguity-sensitive and semantically dependent.
For example, committing to a critical-path witness presupposes commitments that
fix the relative placement of competing spans.

\subsection{Resolving Determinations}

Let $B$ and $C$ be sibling spans whose relative order is not determined by causal
constraints.
Two resolving determinations suffice to illustrate the provenance structure:
\[
  I_B \;\triangleq\;
  \varphi_{\mathrm{align}} \wedge \varphi_{\mathrm{order}(B,C)} \wedge
  \varphi_{\mathrm{cp}=B},
\]
\[
  I_C \;\triangleq\;
  \varphi_{\mathrm{align}} \wedge \varphi_{\mathrm{order}(C,B)} \wedge
  \varphi_{\mathrm{cp}=C}.
\]
Both determinations resolve $\Spec$, but they induce different critical-path
outcomes.

\subsection{Determination Provenance}

Once a determination is fixed, critical-path computation is ordinary.
Let $k_B$ (resp.\ $k_C$) denote the derivational justification that $B$ (resp.\ $C$)
is the longest-duration span in the resolved execution.
Determination provenance records:
\[
  \mathsf{DP}(\mathrm{cp}=B) \;\triangleq\; (I_B,\, k_B),
  \qquad
  \mathsf{DP}(\mathrm{cp}=C) \;\triangleq\; (I_C,\, k_C) .
\]

The provenance makes explicit that the diagnostic conclusion depends on semantic
commitments about clock alignment and span ordering.
Without fixing a determination, the question ``which span is on the critical
path?'' has no determinate answer.

\subsection{Certificates and Parsimony}

Notably, the full trace history is unnecessary to justify the diagnosis.
A sufficient certificate consists of:
\begin{itemize}
  \item the parent--child span structure;
  \item the inequality constraints induced by send/receive events;
  \item the specific alignment and ordering commitments recorded in $I_B$ or
        $I_C$.
\end{itemize}

As in the previous appendices, the certificate acts as a \emph{semantic frontier}:
it freezes exactly those commitments needed to justify the diagnostic observation,
without reconstructing a total execution.
In a totally ordered execution, a timestamp suffices to delimit the past; in a
partially ordered trace, an antichain of events plays the same role, sealing the
relevant downward closure without imposing an artificial linearization.

Determination provenance therefore supports a compressed form of systems lineage:
only the semantic commitments required for the diagnostic conclusion need be
retained.
This refines existing tracing practice by separating observable evidence from
the semantic assumptions used to explain it.

\paragraph{Summary.}
This appendix demonstrates that determination provenance applies naturally to
systems diagnostics.
It extends provenance beyond data derivations to explanations of execution-level
properties, while retaining a minimal, explicit account of the semantic
assumptions required to justify those explanations.

\section{Absence Under Semantic Ambiguity}
\label{app:whynot}
In settings with semantic ambiguity, absence admits distinctions that cannot be
expressed at determination depth~0.

Given a specification $\Spec$ and history $H$, a conclusion may be absent in
several semantically distinct ways:
\begin{enumerate}
  \item It fails under a particular resolved determination but holds under others.
    In the transactional example, tuple $d$ is absent under
    $I_{\mathsf{out}}$ (where $T_5 \prec T_Q$) but present under
    $I_{\mathsf{in}}$ (where $T_Q \prec T_5$).
    Classical why-not provenance under $I_{\mathsf{out}}$ explains the
    derivational failure; determination provenance additionally records that
    the absence is contingent on the ordering commitment.

  \item It fails under \emph{all} resolving determinations of the fixed
    specification $\Spec$.
    In the same example, a tuple $(e)$ that was never inserted fails under
    every serialization.
    This is a robust absence: no ordering commitment can make it appear.
    Determination provenance certifies this by showing that the observation
    lies outside $(\Spec \mid I)(H)$ for every resolving $I$.

  \item It holds under some determinations of a \emph{different} specification
    $\Spec'$ but under none of the determinations admitted by $\Spec$ itself.
    For instance, a write-skew outcome visible under snapshot isolation may be
    absent under every serializable determination.
    This is an inter-specification absence, handled via the preservation
    analysis of Section~\ref{sec:robustness} rather than by determination
    provenance alone.
\end{enumerate}

\noindent
Classical why-not provenance collapses these cases, because it conditions on a
single resolved semantics and a fixed specification.
Determination provenance makes the first two cases explicit by recording the
semantic commitments under which observations hold, explaining absence in terms
of the non-existence (or variability) of supporting determinations rather than
in terms of blocked derivations.
\bibliographystyle{ACM-Reference-Format}
\bibliography{references}

\end{document}
